\documentclass[11pt,a5paper,twoside,openright,headings=small]{scrbook}
\usepackage[T1]{fontenc}
\usepackage{lmodern}
\usepackage[a5paper,inner=2.2cm,outer=1.8cm,top=2.2cm,bottom=2.2cm]{geometry}
\usepackage{microtype}
\usepackage{amsmath,amssymb}
\usepackage{graphicx}
\usepackage[backend=biber,style=authoryear]{biblatex}
\addbibresource{refs.bib}
\usepackage[hidelinks]{hyperref}

\title{The Title of the Book}
\subtitle{A Subtitle That Says What It Is About}
\author{Your Name}
\date{\the\year}

\begin{document}

\frontmatter
\maketitle

\begin{titlepage}\thispagestyle{empty}
  \null\vfill
  \small
  Copyright \copyright{} \the\year{} Your Name\par
  First edition. All rights reserved.\par
  Typeset with \LaTeX{} in Lemmata.
\end{titlepage}

\cleardoublepage
\thispagestyle{empty}
\null\vspace{4cm}
\begin{center}\itshape For whoever this is for.\end{center}

\cleardoublepage
\tableofcontents

\chapter*{Preface}
\addcontentsline{toc}{chapter}{Preface}
Say who the book is for, how it is organised and what the reader needs to know first.
Keep it to a page or two.

\mainmatter

\part{Foundations}
\chapter{The First Chapter}\label{ch:first}

\begin{quote}\itshape
An epigraph, if you like one.
\end{quote}

\section{An Opening Section}
Start with a story or a question, not with definitions. Chapters in a book are
usually 15 to 30 pages; each should end on something the reader can carry forward.
A citation looks like this: \textcite{lamport94} describe the system this book is typeset with.

\section{A Section with Maths}
A displayed formula is numbered by chapter:
\begin{equation}
  \sum_{k=1}^{n} k = \frac{n(n+1)}{2}.
  \label{eq:gauss}
\end{equation}
Cross-references work across files, as in Chapter~\ref{ch:second} and equation~\eqref{eq:gauss}.

\subsection{A Subsection}
Use three levels at most. \KOMAScript's \texttt{scrbook} also provides \verb|\paragraph|
for run-in headings.

\chapter{The Second Chapter}\label{ch:second}

\section{Building on Chapter~\ref{ch:first}}
Refer back to earlier material with \verb|\ref|; the chapter number is filled in for you.

\section{Figures}
Put image files in a \texttt{figures} folder and include them
inside a \texttt{figure} environment that has a caption.

\section{What Comes Next}
Close the chapter with a short summary and a pointer to the next one.


\part{Applications}
% Add chapters as \input{chapters/03-name} and create the file in the chapters folder.

\appendix
\chapter{Notation}
A table of symbols or a longer derivation that would interrupt the main text.

\backmatter
\printbibliography[heading=bibintoc]

\end{document}
